\documentclass[10pt,a4paper]{amsart}
\usepackage[margin=19mm]{geometry}
\usepackage{amsmath,amssymb,mathtools}
\usepackage[scr=rsfs,cal=euler]{mathalfa}
\usepackage{xfrac}
\usepackage[unicode,hidelinks,bookmarks=false]{hyperref}
\newcommand{\ie}{i.e.\ }
\newcommand{\cf}{cf.\ }
\newcommand{\resp}{respectively\ }
\newcommand{\Subsection}{\subsection}
\providecommand{\nobreakdash}{\nobreak\hbox{-}}
\setlength{\emergencystretch}{2em}
\multlinegap0pt
\usepackage{amssymb}
\usepackage{mathtools}
\usepackage{tikz}
\usepackage{xcolor}
\usepackage[shortlabels]{enumitem}
\newtheorem{theorem}{Theorem}
\newtheorem{lemma}{Lemma}
\newcommand{\C}{\mathbb{C}}
\renewcommand{\Im}{\operatorname{Im}}
\newcommand{\N}{\mathbb{N}}
\newcommand{\R}{\mathbb{R}}
\renewcommand{\Re}{\operatorname{Re}}
\newcommand{\boC}{\mathcal{C}}
\title{Stability and instability of the quasilinear Gross--Pitaevskii dark solitons}
\author{Erwan Le Quiniou}
\date{Source: arXiv:2402.18316v4 (CC BY 4.0)}
\begin{document}
\maketitle

\noindent\emph{Source and licence.} Stability and instability of the quasilinear Gross--Pitaevskii dark solitons. Version arXiv:2402.18316v4 (CC BY 4.0), distributed by arXiv under the Creative Commons Attribution 4.0 International licence, http://creativecommons.org/licenses/by/4.0/, as recorded in the arXiv licence metadata for this exact version and shown on the abstract page https://arxiv.org/abs/2402.18316v4. Full licence notice: \texttt{licenses/arxiv-2402.18316-CC-BY-4.0.txt}.\par\smallskip

\noindent\textit{Editorial scope.} Counted unit: the article's lemma lem:dequiv (source0.tex lines 787--792). The article's complete original proof of it is delivered with it (source0.tex lines 793--836, one \texttt{proof} environment of 44 lines). No mathematics is written, repaired or re-derived: the statement and the proof are the article's own lines, in the article's own notation, and no retained line is substituted or re-flowed.  Retained uncounted as the source-local context the proof needs: lines 352--355; lines 357--372; lines 367--369.  Nothing was omitted from any retained span, so the package carries no omission record.  The retained lines cite 3 works, whose entries are transcribed verbatim from the article's own bibliography source in the e-print and delivered with short editorial labels recorded in the item's reference mapping.  No retained line contains a figure environment, an image-inclusion command or drawing code, so no image is delivered and the package declares no asset dependency.  Editorial formatting only, in addition: an AMS \texttt{amsart} preamble that restates the article's own declarations and macros used by the retained lines, namely the environments \texttt{theorem}, \texttt{lemma} on separate counters, together with the standard packages those lines need.  The retained environments print in this document as Theorem 1, Lemma 1 in order of appearance, whereas the article prints the same items with the numbering of its own full document.  The complete native source of the article as distributed in the arXiv e-print for this exact version is bundled unchanged as \texttt{sources/155276/source0.tex}.
	\begin{equation}
		\label{TWc}\tag{TW$(c,\kappa)$}
		-icu_{c,\kappa}'+u_{c,\kappa}''+u_{c,\kappa} (1-|u_{c,\kappa}|^2)-\kappa u_{c,\kappa}\left(|u_{c,\kappa}|^2\right)''=0.
	\end{equation}

	\begin{theorem}[Theorem~1, Proposition~3.8 in \cite{delaire2023exotic}]
		\label{thm:classiftw}
		Let $\kappa\in(-\infty,1/2)$ and $c\geq0$.
		\begin{enumerate}
			\item If $c\geq\sqrt{2}$ and $u_{c,\kappa}\in\mathcal{X}(\R)\cap\boC^{2}(\R)$  satisfies \eqref{TWc}, then $u_{c,\kappa}$ is trivial, i.e.\ there exists $\varphi\in\R$ such that 
			$$u_{c,\kappa}(x)=e^{i\varphi},\quad\text{ for all } x\in\R.$$  
			\item If $0\leq c<\sqrt{2}$ so that $(c,\kappa)\in\mathcal{D}$, then there exists a nonconstant function $u_{c,\kappa}\in\mathcal{X}(\R)\cap\boC^{2}(\R)$ called dark soliton that satisfies \eqref{TWc}.
			Moreover, this solution is unique up to translation and phase change, i.e.\
			any other solution $\phi\in\mathcal{X}(\R)\cap\boC^{2}(\R)$ satisfies for some $(x_0,\varphi)\in\R\times\R$ $$\phi(x)=u_{c,\kappa}(x-x_0)e^{i\varphi},\quad\text{ for all }x\in\R.$$ 
			Additionally, the function $(c,\kappa,x)\mapsto u_{c,\kappa}(x)$ is smooth in $\mathcal{D}\times\R$ (where the derivative at $c=0$ is in the sense of Dini) and satisfies the following decay estimate: for every multi-index $\alpha=(\alpha_1,\alpha_2,\alpha_3)\in\N^3$ and $D^\alpha=\partial_c^{\alpha_1}\partial_\kappa^{\alpha_2}\partial_{x}^{\alpha_3}$, there exist $A,C>0$ continuously depending on $(c,\kappa)$ such that
			\begin{equation}\label{eq:decaydark}
				|D^\alpha\partial_xu_{c,\kappa}|+|D^\alpha(1-|u|^2)|\leq Ae^{-C|x|},\quad\text{ for all }x\in\R.
			\end{equation}
			Also $|u_{c,\kappa}|$ is even with $|u_{c,\kappa}|'>0$ in $(0,\infty)$. Finally we have the bounds $1>|u_{c,\kappa}(x)|\geq c/\sqrt{2}$ for all $x\in\R$ with equality at $x=0$, so that $\inf_\R|u_{c,\kappa}|>0$ if and only if $c>0$.
		\end{enumerate} 
	\end{theorem}

			\begin{equation}
				|D^\alpha\partial_xu_{c,\kappa}|+|D^\alpha(1-|u|^2)|\leq Ae^{-C|x|},\quad\text{ for all }x\in\R.
			\end{equation}

	\begin{lemma}\label{lem:dequiv}
		Let $\kappa<1/2$, $c\in(0,\sqrt{2})$ and $u_{c,\kappa}$ be a dark soliton given by Theorem~\ref{thm:classiftw}. Then there exists $K$ and $\varepsilon>0$ depending only on $u_{c,\kappa}$ such that for every $\Psi\in\mathcal{X}(\R)$, if $\Psi$ satisfies $d_\mathcal{X}(\Psi,u_{c,\kappa})<\varepsilon$, then we can write $\Psi=\sqrt{1-\eta}e^{i\theta}$ with $(\eta,\partial_x\theta)\in X$ and we have
		\begin{equation}
			\frac{1}{K}d_{\mathrm{hy}}(\Psi,u_{c,\kappa})\leq d_{\mathcal{X}}(\Psi,u_{c,\kappa})\leq Kd_{\mathrm{hy}}(\Psi,u_{c,\kappa}).
		\end{equation} 
	\end{lemma}

	\begin{proof}
		For clarity's sake, we recall the main arguments of the proof in \cite{chironstab}.
		We remove $\kappa$ in the subscript of $u_{c,\kappa}=u_c$ and write $u_c=\sqrt{1-\eta_c}e^{i\theta_c}$. Since $u_{c}$ does not vanish, taking $\varepsilon>0$ small enough we deduce that for any $\Psi\in\mathcal{X}(\R)$ satisfying $d_\mathcal{X}(\Psi,u_c)\leq\varepsilon$, we have $|||\Psi|-|u_{c}|||_{ L^\infty(\R)}\leq\inf_{x\in\R}|u_c(x)|/2$ so that 
		\begin{equation}\label{eq:closinf}
			\inf_{x\in\R}|\Psi(x)|\geq \inf_{x\in\R}|u_c(x)|/2.
		\end{equation} We can express $\partial_x\Psi$ in terms of $\eta$ and $\theta$,
		\begin{equation}\label{eq:Psiprim}
			\partial_x\Psi=e^{i\theta}\left(\frac{\eta_x}{2\sqrt{1-\eta}}+i\sqrt{1-\eta}\theta_x\right),
		\end{equation} and similarly for $u_c$.
		We also need to bound $||\Psi||_{ L^\infty(\R)}$ in terms of $\varepsilon$, $||u_c||_{ L^\infty(\R)}$ and $||\partial_xu_c||_{ L^2(\R)}$. To do so, notice that for any constants $y<x$, we have
		\begin{equation}\label{eq:morrey}
			|\Psi(x)|\leq|\Psi(y)|+\int_{y}^x|\partial_x\Psi|\leq|u_c(y)|+||\Psi(y)|-|u_c(y)||+\sqrt{x-y}||\partial_x\Psi||_{ L^2(\R)},
		\end{equation} where we used Cauchy--Schwarz inequality to bound the integral. Taking $0<x-y<1$ then integrating \eqref{eq:morrey} with respect to $y$ from $x-1$ to $x$ we get
		$|\Psi(x)|\leq |u_c|+\int_{\R}||u_c|-|\Psi||+||\partial_x\Psi||_{ L^2(\R)},$
		so that
		\begin{equation}\label{eq:inftybound}
			||\Psi||_{ L^\infty(\R)}\leq K_0(u_c,\varepsilon).
		\end{equation} 
		We show that $d_{\mathrm{hy}}(\Psi,u_{c})\leq Kd_\mathcal{X}(\Psi,u_c).$ 
		First notice that $|\Psi|^2\theta_x=\Im(\Psi_x\bar{\Psi})$, thus
		$$\partial_x\theta-\partial_x\theta_c=\Im(\partial_xu_c\bar{u_c})\left(\frac{1}{|\Psi|^2}-\frac{1}{|u_c|^2}\right)+\frac{1}{|\Psi|^2}(\Im((\partial_x\Psi-\partial_xu_c)\bar{\Psi})-\Im(\partial_xu_c(\bar{u_c}-\bar{\Psi}))),$$
		
		\begin{flalign}\label{eq:boundthetax}
			\text{hence, }\hspace{1.3cm} ||\partial_x\theta-\partial_x\theta_c||_{ L^2(\R)}\leq& \frac{\Im(\partial_xu_c\bar{u_c})(||\Psi||_{ L^\infty(\R)}+||u_c||_{ L^\infty(\R)})}{\inf_{\R}|\Psi|^2\inf_\R|u_c|^2}|||\Psi|-|u_c|||_{ L^2(\R)}&&\notag\\
			&+\frac{1}{\inf_\R|\Psi|^2}\left(||\Psi||_{ L^\infty(\R)}||\partial_x\Psi-\partial_xu||_{ L^2(\R)}+||\partial_xu_c(\bar{u_c}-\bar{\Psi})||_{ L^2(\R)}\right).
		\end{flalign}
		Since $\Psi(x)-u_c(x)=\Psi(0)-u_c(0)+\int_{0}^x\partial_x\Psi-\partial_xu_c$ for all $x\in\R$, using the Cauchy--Schwarz inequality, we get
		\begin{equation}
			|\partial_xu_c(x)||\Psi(x)-u_c(x)|\leq|\partial_xu(x)||\Psi(0)-u_c(0)|+|\partial_{x}u_c(x)\sqrt{|x|}|||\partial_x\Psi-\partial_xu_c||_{ L^2(\R)},\quad\text{ for all }x\in\R.
		\end{equation}
		From the decay estimate~\eqref{eq:decaydark} satisfied by $\partial_xu_c$, we can bound the last term in \eqref{eq:boundthetax}, therefore using also \eqref{eq:closinf} and \eqref{eq:inftybound}, we deduce that there exists $K_1=K_1(\varepsilon,u_c)$ such that 
		\begin{equation}\label{eq:equivtheta}
			||\partial_x\theta-\partial_x\theta_c||\leq K_1d_{\mathrm{hy}}(\Psi,u_c).	
		\end{equation}
		Similarly, we have $2\Re(\Psi_x\bar\Psi)=\eta_x$, thus 
		\begin{align}
			||\partial_x\eta-\partial_x\eta_c||_{ L^2(\R)}&\leq2||\Psi||_{ L^\infty(\R)}||\partial_x\Psi-\partial_xu_c||_{ L^2(\R)}+2||\partial_x{u_c}(\bar{u_c}-\bar\Psi)||_{ L^2(\R)},\notag\\
			&\leq2K_0d_{\mathrm{hy}}(\Psi,u_c)+2||\partial_xu_c||_{ L^2(\R)}||\Psi(0)|-|u_c(0)||+||\sqrt{x}\partial_xu_c||_{ L^2(\R)}|| \partial_x\Psi-\partial_xu_c||_{ L^2(\R)}\notag,\\
			&\leq K_2 d_{\mathrm{hy}}(\Psi,u_c).
		\end{align}
		Then, we also have $||\eta-\eta_c||_{ L^2(\R)}=||(|\Psi|+|u_c|)(|\Psi|-|u_c|)||_{ L^2(\R)}\leq(||\Psi||_{ L^\infty(\R)}+||u_c||_{ L^\infty(\R)})|||\Psi|-|u_c|||_{ L^2(\R)}$.
		We bound the last term $|\mathrm{Arg}(\Psi(0)/u_c(0))|$. Letting $\mathrm{Log}$ be the principal determination of the logarithm defined in $\C\backslash(-\infty,0]$ and taking $\varepsilon>0$ small enough, we get $$|\mathrm{Arg}(\Psi(0)/u_c(0))|=|\Im (\mathrm{Log}(\Psi(0)/u_c(0)))|\leq \sup_{|z-1|\leq\varepsilon/u_c(0)}|1/z|\left|\frac{\Psi(0)-u_c(0))}{u_c(0)}\right|,$$ which ends to prove that $d_{\mathrm{hy}}(\Psi,u_c)\leq K d_{\mathcal{X}}(\Psi,u_c)$.
		Ideas along the same lines as in Lemma~10 in \cite{chiron-stability} show that $d_{\mathcal{X}}(\Psi,u_c)\leq K d_{\mathrm{hy}}(\Psi,u_c)$.
	\end{proof}

\begin{thebibliography}{99}
\bibitem[chiron]{chiron-stability} D.~Chiron. \newblock Stability and instability for subsonic traveling waves of the nonlinear {S}chr\"{o}dinger equation in dimension one. \newblock {\em Anal. PDE}, 6(6):1327--1420, 2013.
\bibitem[chiron]{chironstab} D.~Chiron. \newblock Stability and instability for subsonic traveling waves of the nonlinear {S}chr\"odinger equation in dimension one. \newblock {\em Anal. PDE}, 6(6):1327--1420, 2013.
\bibitem[delair]{delaire2023exotic} A.~de~Laire and E.~L. Quiniou. \newblock Exotic traveling waves for a quasilinear {S}chr\"odinger equation with nonzero background, 2023. \newblock Preprint arXiv:2311.08918.
\end{thebibliography}
\end{document}
